\section{Discussion}

\subsection{What the contradictions are, and are not}

Taken together, the results support a single claim: the apparent inter-study contradictions
in aluminum--water hydrolysis are \emph{predominantly methodological} (though this attribution is observational and
may partly reflect unmeasured study-level factors). Naive pooling
overstates predictability by $0.62$--$0.85$ in $R^2$; the physico-chemical regime label
explains $\approx\!2\%$ of between-study yield variance while methodological covariates
explain $\approx\!72\%$; the flagship particle-size contradiction dissolves. The practical
implication is direct: a hydrogen yield reported without its measurement protocol is only
weakly comparable across laboratories, and meta-analyses that pool raw yields inherit
methodological variance as if it were chemistry.

These two analyses are complementary, not competing. The between-study variance
decomposition (T\_B) is a study-level regression ($n\!=\!31$ aggregated study rows) that
attributes study-mean yield differences to study-level covariates. The mixed-effects
meta-regression (H2) operates on row-level data ($n\!=\!299$) with a random intercept
per study: that intercept absorbs all unmeasured between-study heterogeneity, including
methodology and study identity together, leaving regime as the sole tested fixed effect.
The variance decomposition answers ``what covariates explain why studies differ?''; the
mixed-effects model answers ``does regime moderate yield, after attributing all other
between-study variation to a random effect?'' Both reach the same conclusion: the
physico-chemical regime label carries weak explanatory power relative to study-level
methodological choices.

Naming the specific apparent contradictions helps calibrate the claim. For activation
energy, the most striking contrast in the liquid-metal regime is between
\citeauthor{ilyukhina2012}~\citep{ilyukhina2012} ($E_\mathrm{a}\!=\!8.5$\,kJ\,mol$^{-1}$,
Ga--In eutectic, peak-rate proxy) and \citeauthor{dudoladov2016}~\citep{dudoladov2016}
($E_\mathrm{a}\!=\!58.0$\,kJ\,mol$^{-1}$, Ga in KOH, tabulated rate): both are
liquid-metal-activated systems, yet they differ by nearly sevenfold. Our analysis does
not resolve whether this reflects different Ga/In loading, different activation proxies
(peak rate vs.\ initial rate), or different surface geometries; it demonstrates that the
current ``liquid\_metal\_activated'' label is too coarse to resolve this.
For particle size, cross-study comparisons are misleading because studies compare
incompatible size windows: one laboratory's ``fine'' particles are another's standard
size, and different regimes are mixed when averaging across studies. Within each of the
five studies that vary particle size, the effect is consistent and negative (smaller
$\rightarrow$ higher yield; Spearman $\rho\!=\!-0.29$ to $-0.87$ in
\citep{zhang2024,prabu2021,wen2018,porciuncula2012,deng2010}). The apparent cross-study flip is
thus an artifact of comparing different experimental windows, not a genuine physics
disagreement.

\subsection{Why pooling fails: the cross-study evidence in context}

The scale of the optimism gap ($0.62$--$0.85$) is consistent with what has been found
when grouped cross-validation is applied to other domains with hidden group structure
\citep{bernett2024,john2025}. The closest published analogues for hydrogen generation
are the interlaboratory round-robin comparisons for MgH$_2$ storage \citep{wulf2021} and
biohydrogen fermentation \citep{hoque2018}: in both cases, identical ``conditions'' in
different laboratories produced yields spanning roughly twofold, with the spread traced
primarily to measurement and protocol choices rather than to genuine material differences.
The present result, that temperature control alone explains $33\%$ of between-study
yield variance in aluminum hydrolysis, places this field in the same category. The
leakage-controlled benchmark we release is designed to let future modeling efforts detect
the same failure mode: any model that achieves positive grouped $R^2$ on the provided
corpus under the study-grouped split would represent genuine cross-study generalization
beyond what the current literature supports.

\subsection{Physics is not absent}

``Predominantly methodological'' is not ``purely methodological.'' The regime label is
coarse; finer, particle-size-dependent kinetics (tunable between nano- and
micro-aluminium powders) genuinely shape the reaction \citep{saceleanu2019}, and such
physics operates \emph{within} studies, below the resolution of a five-level label and
partly absorbed into the methodological covariates with which it co-varies. Incremental
regime variance ($+0.15$ after controlling for method) is nonzero, though its precise
magnitude cannot be established reliably at $n\!=\!31$ studies.
Our attribution is also observational: because methodology is largely a study-level
attribute, the temperature-control association could be partly confounded with other
unmeasured study-level choices (principal investigator laboratory practice, instrument
calibration, student experience). We state the finding associationally and quantify,
rather than assert, its robustness.

\subsection{A minimum-information reporting standard}

The remedy is standardization. Table~\ref{tab:mireport} lists the minimum information
needed to interpret a hydrogen yield across laboratories. The ordering reflects the
between-study variance contribution each field carried in our analysis: temperature
control is listed first because it was the single largest source ($R^2\!=\!0.33$);
measurement method and water type follow because they contribute smaller but nonzero
shares. The remaining fields (alkali type and concentration, particle size, activation
route, yield basis, rate definition) carried detectable but weaker or less-covered
signal here, yet they are practically essential for reproducibility. Provenance (DOI,
condition count, replicate count) is necessary to build cumulative, poolable evidence.

This template is modelled on reporting standards that successfully improved poolability
in related fields: round-robin-derived checklists for MgH$_2$ \citep{wulf2021} and
biohydrogen \citep{hoque2018} both show that standardization shrinks interlaboratory
spread rather than eliminating it. We expect the same here: full compliance would not
remove between-study variance, but it would expose the remaining variance as physics
rather than protocol.

\begin{table}[t]\centering\small
\caption{Proposed minimum-information reporting standard for aluminum--water hydrolysis
hydrogen-yield studies. Items are ordered by the between-study yield variance each
carried here (Section~\ref{sec:variance_decomp}); temperature control is the largest
single source ($R^2\!=\!0.33$).}
\label{tab:mireport}
\begin{tabular}{lll}
\toprule
Field & Report & Between-study $R^2$ here \\
\midrule
Temperature control  & method (isothermal bath / self-heating / uncontrolled) + setpoint & 0.33 \\
Measurement method   & water displacement / pressure / GC / mass-flow & 0.07 \\
Water type \& purity & DI / tap / alkaline / saline, with conductivity or grade & 0.04 \\
Quality / control    & replication count; value reported vs.\ derived & 0.04 \\
Reaction vessel      & open or closed; insulated? & 0.02 \\
Alkali               & type (NaOH/KOH/none) + concentration (mol\,L$^{-1}$) & --- \\
Particle size        & $D_{50}$ ($\mu$m) + sizing method & --- \\
Activation route     & pure / melt-alloyed / ball-milled / liquid-metal / waste & --- \\
Yield basis          & theoretical reference (mL\,g$^{-1}$); \% vs.\ volumetric & --- \\
Rate definition      & initial / maximum / average; $t_{50}$/$t_{80}$ if rate is claimed & --- \\
Provenance           & study DOI, number of conditions, replicate count & --- \\
\bottomrule
\end{tabular}
\end{table}

\subsection{Limitations}
\label{sec:limitations}

\paragraph{Sample size and power.}
The analysis rests on $n\!\approx\!31$ studies ($3$--$9$ per regime). Null results
(no regime moderation of yield, H2; no regime structure in $E_\mathrm{a}$, H3;
heterogeneous within-study predictability, T\_A) are reported as ``no detectable effect
at this sample size,'' not as proof of absence. Detection of small effects ($R^2\!<\!0.10$ for regime) would
likely require a corpus substantially larger than $n\!=\!31$ studies.
The between-study variance decomposition
(T\_B) has cluster-bootstrap confidence intervals that are wide at this $n$; the lower
bound on method-joint $R^2$ ($0.44$) excludes zero but the upper bound ($1.00$) is
uninformative. Two regimes (\texttt{waste\_al}, \texttt{liquid\_metal\_activated}) have
fewer than $40$ yield rows each; effect estimates for these classes are exploratory.

\paragraph{Associational attribution.}
All findings are associational, not causal. Because measurement apparatus and temperature-control
method are study-level attributes, their variance shares ($43\%$ and $8\%$ respectively)
could partly reflect unmeasured study-level differences (equipment, calibration, laboratory
practice) that co-vary with protocol choice. We quantify the incremental contributions
(measurement method beyond regime: $+0.51$; temperature control beyond regime: $+0.11$;
regime beyond both: $+0.04$) but do not attribute a causal mechanism.

\paragraph{Single material system and open-access bias.}
The dataset covers only aluminum--water hydrolysis; the leakage finding and the
methodological-heterogeneity finding may not transfer to other reaction systems, though
analogues exist in MgH$_2$ and biohydrogen \citep{wulf2021,hoque2018}. The 31 studies
are the open-access-reachable subset of the screened literature, not a complete or
random sample. Open-access status may correlate with publication venue prestige and
reporting completeness, but there is no clear mechanistic reason why it should
systematically predict which temperature-control method or measurement protocol a
laboratory adopts; the null expectation is therefore that OA-only retrieval does not
directionally bias the primary methodology-versus-regime finding. The direction of
residual bias cannot be verified from the current corpus; external validation using a
broader sample would resolve this.

\paragraph{Unvalidated reporting standard.}
The proposed minimum-information checklist (Table~\ref{tab:mireport}) has not been
prospectively validated: we do not yet know how much of the residual between-study
variance it would remove in practice. It is a hypothesis about what would help, grounded
in our variance decomposition. Prospective validation would require a multi-laboratory
round-robin trial in which independent groups run the same experiment under the
standard's reporting protocol \citep[analogous to those for MgH$_2$ and
biohydrogen;][]{wulf2021,hoque2018}; this is beyond the scope of the present work
but is a logical next step.

\paragraph{Single-author and AI-assisted workflow.}
All data extraction, curation, analysis, and writing were performed by a single author
with AI assistance (see Declaration of generative AI). Single-author extraction provides
intra-rater consistency evidence only, not inter-rater reliability: the double-extraction
check ($52/52$ exact, $52/315$ rows, $16\%$) repeated extraction by the same author on a
subset of rows, confirming self-consistency but not guarding against systematic extractor
bias shared across both passes. The remaining $84\%$ of rows ($263/315$) have no
independent extractor check. A post-hoc independent-agent IRR pass (Section~\ref{sec:qc})
identified three coding errors in the study-level metadata
(\texttt{temperature\_control} and \texttt{measurement\_method})
that, after correction, meaningfully changed the variance decomposition numbers
(temperature control: $0.332\to0.081$; measurement apparatus: nominally largest predictor,
$0.425$). This underscores that agent-assisted IRR can surface coding errors even when
the original pass was author-reviewed, and that the variance decomposition is sensitive
to the accuracy of study-level metadata labels. The IRR pass covers an accessible
subset ($8/31$ openly reachable studies); systematic biases in the paywalled studies
($n\!=\!23$) cannot be ruled out. The agent-based IRR further shares a base model
(Anthropic Claude), so model-level systematic errors are not excluded; this is not
equivalent to human dual-coding. Readers should treat the results as a first
systematic synthesis requiring external validation rather than as an established consensus.
The curated data and analysis pipeline are released openly so that independent groups can
audit extraction quality and extend the corpus.

\paragraph{Particle-size finding scope.}
The resolution of the particle-size contradiction rests on $n\!=\!5$ studies
($59/315$ rows, $19\%$) that systematically vary particle size. Two regimes
(\texttt{al\_alloy}, \texttt{mechanically\_activated}) have zero particle-size entries.
The $5/5$ within-study consistency is a strong qualitative signal, but the $26$ studies
that did not report particle size may differ systematically from those that did.
Whether the consistently-negative within-study effect generalises beyond the reporting
subset cannot be established from the current dataset.

\paragraph{Leakage gap as a conflated signal.}
The optimism gap ($0.62$--$0.85$ for tree models) reflects two distinct phenomena:
between-study methodological heterogeneity and out-of-distribution extrapolation.
If the 31 studies explore non-overlapping regions of physicochemical parameter
space (plausible given five regimes with different activation methods, particle sizes,
and alkali types), a model trained on some studies will fail on others even under
identical measurement protocols, because the held-out study's conditions are out of
distribution. We cannot currently separate these contributions quantitatively; the
$72\%$ methodology attribution addresses the heterogeneity component but does not
isolate it from extrapolation effects.

\paragraph{Mixed-effects model specification.}
The mixed-effects meta-regression (H2) assumes homoscedastic within-study residual
variance. However, the within-study LOSO analysis (T\_A) shows wide heterogeneity in
within-study predictability ($\eta^2_{\text{study}}\!=\!0.496$). If residual variance
differs substantially across studies, REML estimates may be inefficient and fixed-effect
standard errors biased, making the null regime effect harder to interpret. Influence
diagnostics and heteroscedastic alternatives were not performed; the H2 null should be
replicated in a larger sample.

\paragraph{Laboratory clustering.}
The study-grouped cross-validation treats each DOI as an independent unit. Multiple
papers from the same laboratory or principal investigator may share unmeasured
similarities (equipment, protocols, calibration) that DOI-level grouping does not
capture. If studies cluster by laboratory, the observed $R^2_{\text{grouped}}$ may
still overestimate true cross-laboratory generalization. A lab-grouped cross-validation
was not performed; future work should construct a study-to-laboratory mapping and
test whether the optimism gap widens under stricter grouping.
